\documentclass[11pt]{article}

\usepackage[margin=1in]{geometry}
\usepackage{booktabs}
\usepackage{hyperref}
\usepackage{natbib}
\usepackage{graphicx}
\usepackage{array}
\usepackage{longtable}
\usepackage{microtype}

\title{FreightBidBench: A Public-Calibrated Benchmark for Real-Time Truckload Bid Acceptance Under Operational Feasibility Constraints}
\author{Aswin Sekar}
\date{Draft v0.2.0 -- May 6, 2026}

\begin{document}
\maketitle

\begin{abstract}
Truckload carriers and brokers must make accept/reject decisions under tight
latency constraints, uncertain future demand, and operational feasibility
constraints. Existing academic formulations often study dynamic fleet
management or stochastic vehicle routing, but reproducible public benchmarks
for real-time truckload bid evaluation remain limited. We introduce
FreightBidBench, a public-calibrated synthetic benchmark built from FAF
freight-flow structure and USDA AMS truck-rate reports. FreightBidBench
evaluates policies in closed loop: accepted loads move trucks through the
network, consume driver availability, and affect future opportunity cost.
Version 0.2 adds individual truck state, pickup reach time, pickup and delivery
windows, simplified hours-of-service clocks, and stochastic shipper/receiver
yard delays. The benchmark reports profit, latency, rollout-call share,
infeasible accept attempts, HOS rest hours, deadhead miles, yard-delay hours,
and profit retention relative to a finite rollout teacher. We provide a
dependency-free reference runner, manifest format, baseline policies, and
latency-profit frontier outputs.
\end{abstract}

\section{Introduction}

Real-time truckload bid acceptance is a high-pressure sequential decision
problem. A carrier can accept a load that looks profitable immediately but
sends a truck into a weak future market, violates pickup feasibility, or
consumes driver hours needed for a better future load. Conversely, rejecting a
low-margin load can be costly when the destination improves fleet positioning.

The main obstacle to reproducible research in this setting is not only model
choice. It is the lack of a public benchmark that combines stochastic demand,
closed-loop fleet state, bid-evaluation latency, and operational feasibility.
Private tender data is rarely shareable. Production dispatch systems are
complex and unavailable to most researchers.

FreightBidBench addresses this gap with a public-calibrated benchmark rather
than a claim of production-grade simulation. The purpose is to create a common
testbed where researchers can compare policies on the same stochastic freight
environment and report the same latency-profit metrics.

\section{Contributions}

This paper makes five contributions.

\begin{enumerate}
  \item We introduce FreightBidBench, a reproducible public-calibrated
  benchmark for real-time truckload bid acceptance.
  \item We define a closed-loop synthetic freight environment calibrated from
  public FAF and USDA-derived inputs.
  \item We add a v0.2 feasibility layer with individual trucks, pickup reach
  time, pickup and delivery appointment windows, simplified HOS clocks, and
  stochastic yard delays.
  \item We provide reference policies: myopic margin, bid-price heuristic,
  linear rollout-label surrogate, selective surrogate/rollout cascade, and
  finite rollout teacher.
  \item We define benchmark reporting rules: seed protocol, manifest,
  latency-profit frontier, feasibility metrics, and confidence intervals.
\end{enumerate}

\section{Related Work}

\paragraph{Dynamic truckload acceptance and fleet management.}
Dynamic truckload routing and load acceptance has long been recognized as an
online, state-dependent decision problem. Kim, Mahmassani, and Jaillet study
dynamic load acceptance and rejection decisions for large-fleet truckload
operations with time windows and priority demand \citep{kim2004dynamic}.
Simao et al. show how approximate dynamic programming can estimate marginal
driver values in large-scale truckload fleet management \citep{simao2009adp}.
FreightBidBench builds on this tradition but changes the artifact: instead of
proposing one dispatch method, it provides a benchmark where future-value
policies, rollout teachers, and fast approximations can be compared under the
same stochastic seeds and latency measurements.

\paragraph{Dynamic and stochastic vehicle routing.}
Surveys by Pillac et al. and Ritzinger et al. describe the broader dynamic and
stochastic vehicle-routing literature, including the role of information
evolution and the distinction between online computation and offline policy
construction \citep{pillac2013review,ritzinger2016survey}. Ulmer et al.
combine offline value function approximation with online rollout for dynamic
routing with stochastic requests \citep{ulmer2019offline}, while Ulmer and
Thomas study value-function approximations for dynamic customer acceptance in
delivery routing \citep{ulmer2020meso}. These papers motivate
FreightBidBench's finite rollout teacher and surrogate baselines. The
benchmark gap is that truckload bid evaluation also requires public
calibration, fleet repositioning, operational feasibility, and latency-profit
reporting.

\paragraph{Learning-augmented optimization.}
Recent work uses learned models to speed up routing and stochastic
optimization. Nazari et al. and Kool et al. show that neural policies can
learn routing heuristics with fast inference after training
\citep{nazari2018rlvrp,kool2019attention}. Patel et al. approximate expensive
two-stage stochastic-programming value functions with neural surrogates
\citep{patel2022neur2sp}. Derrow-Pinion et al. demonstrate graph neural
networks for large-scale transportation ETA prediction
\citep{derrowpinion2021eta}. These works support future FreightBidBench
submissions, but the benchmark contribution is not a neural architecture. It is
a reproducible environment and reporting protocol for closed-loop truckload bid
evaluation.

\paragraph{Public freight data and feasibility rules.}
FreightBidBench uses the Freight Analysis Framework as a public freight-flow
backbone \citep{btsfaf5} and USDA AMS Specialty Crops National Truck Rate
Reports for refrigerated truck-rate and availability signals
\citep{usdafvwtrk}. The v0.2 feasibility layer uses a simplified
property-carrying HOS abstraction based on FMCSA's 11-hour driving limit,
14-hour driving window, and 10-hour off-duty reset summary \citep{fmcsa_hos}.
These data sources make the benchmark reproducible while also defining its
limits: the generated tenders are synthetic and public-calibrated, not private
carrier histories.

\section{Benchmark Design}

\subsection{Decision Problem}

At each decision epoch, the environment tenders a candidate load. A policy
chooses accept or reject. If accepted, the benchmark attempts to assign a
feasible truck. If no feasible truck exists, the accepted load creates a
no-truck or infeasible outcome. If feasible, the truck moves to the destination
and becomes available after pickup, linehaul, delivery, yard delays, and
required HOS rest.

\subsection{State}

The v0.2 state includes the candidate load origin/destination, load price,
direct cost, lane distance, public-calibrated lane intensity, state-level
future-value features, individual truck locations, truck availability times,
simplified HOS drive and duty clocks, pickup and delivery appointment windows,
pickup deadhead time, and yard-delay realizations.

\subsection{Reward}

For accepted feasible loads, profit is
\[
  \mathrm{profit} = \mathrm{price} - \mathrm{direct\ cost}
  - \mathrm{pickup\ deadhead\ cost} - \mathrm{yard\ delay\ cost}.
\]
Rejected loads receive zero immediate reward and preserve fleet state.

\section{Public Calibration}

FreightBidBench uses FAF state-level freight flows to seed
origin-destination intensity and regional imbalance. It uses USDA AMS FVWTRK
truck-rate reports to seed refrigerated truckload rates and availability
categories. The benchmark does not claim that generated loads are production
tenders. Instead, public data anchors broad flow, rate, and scarcity structure
so the benchmark is reproducible and inspectable.

\begin{table}[t]
\centering
\caption{Public calibration inputs used in FreightBidBench v0.2.}
\begin{tabular}{>{\raggedright\arraybackslash}p{0.17\linewidth}
                >{\raggedright\arraybackslash}p{0.23\linewidth}
                >{\raggedright\arraybackslash}p{0.24\linewidth}
                >{\raggedright\arraybackslash}p{0.24\linewidth}}
\toprule
Source & Benchmark use & Fields used & Main limitation \\
\midrule
FAF5 state freight flows & OD intensity and regional imbalance & Tons,
ton-miles, origin state, destination state, mode & Annual aggregate, not
tender-level \\
USDA AMS FVWTRK & Reefer rate and scarcity seed & Spot truckload rates,
destination cities, truck availability categories & Narrow equipment and
commodity coverage \\
FMCSA HOS summary & Simplified feasibility clock & 11-hour drive, 14-hour
duty, 10-hour reset & Simplified; omits split sleeper, recap, and home-time \\
\bottomrule
\end{tabular}
\label{tab:calibration}
\end{table}

\section{Operational Feasibility Layer}

Version 0.2 adds a first feasibility layer: individual truck state, pickup
reach, pickup and delivery appointment windows, simplified HOS clocks, and
stochastic yard delays. This layer is intentionally simpler than a full dispatch
simulator. It excludes road closures, route-level traffic, split sleeper rules,
driver home time, team drivers, and maintenance.

\section{Reference Policies}

FreightBidBench includes five reference policy families. The myopic margin
policy accepts when immediate margin is non-negative. The bid-price heuristic
accepts when immediate margin plus destination-origin future-value delta clears
zero. The linear surrogate trains a dependency-free ridge linear model on
offline rollout incremental-value labels. The selective cascade uses the
surrogate for most loads and escalates to rollout near the decision boundary or
when truck supply is scarce. The finite rollout teacher evaluates accept and
reject branches using common-random-number stochastic future simulations and a
bid-price base policy. The teacher is a stochastic benchmark, not an oracle.

\section{Metrics}

Primary metrics are closed-loop profit, profit retention versus rollout
teacher, mean and p95 decision latency, and rollout-call share for cascades.
Feasibility metrics are infeasible accept attempts, pickup-window misses,
delivery-window misses, no-truck outcomes, deadhead miles, HOS rest hours, and
yard-delay hours. Statistical reporting includes seed count, train/eval seed
pairs, mean, sample standard deviation, 95 percent confidence intervals, and
the full benchmark manifest.

\section{Experimental Protocol}

FreightBidBench v0.2 defines three scenarios: \texttt{mild}, \texttt{tight},
and \texttt{scarce}. The \texttt{mild} regime is intentionally easy so that
myopic policies can be competitive. The \texttt{tight} and \texttt{scarce}
regimes increase the value of future-aware decisions.

\begin{table}[t]
\centering
\caption{Benchmark presets.}
\begin{tabular}{lll}
\toprule
Preset & Use & Configuration \\
\midrule
\texttt{smoke} & Correctness check & One seed pair, tight scenario, short horizon sample \\
\texttt{standard} & Local development result & Three seed pairs across all scenarios \\
\texttt{paper} & Main paper table & Ten seed pairs across all scenarios \\
\bottomrule
\end{tabular}
\label{tab:presets}
\end{table}

The standard run can be reproduced with:
\begin{verbatim}
python3 scripts/run_freightbidbench.py --preset standard \
  --output-dir benchmark_runs/standard_v02
\end{verbatim}

\section{Results}

The current v0.2 standard run used three train/eval seed pairs across the
\texttt{mild}, \texttt{tight}, and \texttt{scarce} scenarios. It used 600
rollout-label decisions per train/eval stream and evaluated the full 72-hour
online horizon. The run wrote 99 seed-level policy rows, 9 static-fit rows, 33
aggregate policy rows, and 21 cascade-frontier rows. Total runtime was 1,224.07
seconds.

\subsection{Offline Label Fit}

The simple linear surrogate is intentionally weak under the v0.2 feasibility
layer. Held-out accept/reject agreement was 58.8 percent in \texttt{mild},
60.9 percent in \texttt{tight}, and 68.3 percent in \texttt{scarce}. This is
useful for the benchmark paper: the reference surrogate is a baseline, not the
contribution.

\subsection{Policy Results}

\begin{table}[t]
\centering
\caption{FreightBidBench v0.2 standard policy results. Profit, latency, and
feasibility metrics are seed means over three train/eval seed pairs. Retention
is measured against the finite rollout teacher in the same scenario.}
\small
\begin{tabular}{llrrrrr}
\toprule
Scenario & Policy & Profit & Retention & Latency ms & Infeasible & HOS h \\
\midrule
\texttt{mild} & Myopic & \$1,082,891 & 101.7\% & 0.000 & 323.7 & 2,053 \\
\texttt{mild} & Bid price & \$1,083,424 & 101.8\% & 0.001 & 312.0 & 2,050 \\
\texttt{mild} & Linear surrogate & \$917,369 & 86.3\% & 0.004 & 112.3 & 1,540 \\
\texttt{mild} & Cascade +/- \$500 & \$976,891 & 91.9\% & 15.394 & 76.3 & 1,627 \\
\texttt{mild} & Rollout teacher & \$1,064,003 & 100.0\% & 43.857 & 0.0 & 1,840 \\
\texttt{tight} & Myopic & \$866,894 & 92.0\% & 0.000 & 546.7 & 1,680 \\
\texttt{tight} & Bid price & \$866,894 & 92.0\% & 0.001 & 533.7 & 1,680 \\
\texttt{tight} & Linear surrogate & \$739,135 & 78.5\% & 0.004 & 87.3 & 1,217 \\
\texttt{tight} & Cascade +/- \$500 & \$812,794 & 86.3\% & 10.601 & 34.3 & 1,363 \\
\texttt{tight} & Rollout teacher & \$942,219 & 100.0\% & 27.373 & 0.0 & 1,753 \\
\texttt{scarce} & Myopic & \$718,085 & 94.8\% & 0.000 & 754.3 & 1,377 \\
\texttt{scarce} & Bid price & \$718,085 & 94.8\% & 0.001 & 740.0 & 1,377 \\
\texttt{scarce} & Linear surrogate & \$487,441 & 64.4\% & 0.004 & 40.3 & 813 \\
\texttt{scarce} & Cascade +/- \$500 & \$577,343 & 76.3\% & 7.087 & 0.0 & 950 \\
\texttt{scarce} & Rollout teacher & \$757,682 & 100.0\% & 16.751 & 0.0 & 1,403 \\
\bottomrule
\end{tabular}
\label{tab:results}
\end{table}

Two findings matter for the benchmark paper. First, feasibility metrics are not
secondary bookkeeping: myopic and bid-price policies create hundreds of
infeasible accept attempts in the v0.2 standard run. Second, rollout is much
slower under feasibility, with mean latency from 16.751 ms to 43.857 ms across
scenarios. This creates a clear latency-quality benchmark for future methods.

\subsection{Cascade Frontier}

The cascade frontier shows the tradeoff between rollout-call share and profit
retention. At the widest tested band, the cascade reached 97.1 percent
retention in \texttt{mild}, 91.8 percent in \texttt{tight}, and 90.5 percent
in \texttt{scarce}. The result is not a final method claim; it shows that the
benchmark can expose meaningful frontier differences across scenarios.

\begin{figure}[t]
\centering
\fbox{\begin{minipage}{0.88\linewidth}
\centering
Generated SVG: \texttt{benchmark\_runs/standard\_v02/figures/profit\_retention\_by\_policy.svg}
\end{minipage}}
\caption{Profit retention by policy.}
\label{fig:profit-retention}
\end{figure}

\begin{figure}[t]
\centering
\fbox{\begin{minipage}{0.88\linewidth}
\centering
Generated SVG: \texttt{benchmark\_runs/standard\_v02/figures/latency\_profit\_frontier.svg}
\end{minipage}}
\caption{Latency-profit frontier.}
\label{fig:latency-frontier}
\end{figure}

\section{Discussion}

The central benchmark claim is that feasibility materially changes the
accept/reject problem. A policy that looks attractive under aggregate
state-level availability can create pickup-window misses, excess HOS rest, or
yard-delay exposure once individual truck feasibility is considered.

The latency result is also important: operational feasibility increases rollout
runtime. This makes FreightBidBench useful for studying selective evaluation,
surrogates, and other learning-augmented acceleration methods.

\section{Limitations}

FreightBidBench v0.2 is not a production dispatch simulator. It uses synthetic
loads, public calibration, simplified HOS clocks, state-level geography, and
reference Python latency. It does not model road closures, traffic, weather,
team drivers, split sleeper rules, maintenance, customer-specific facilities,
or private tender distributions.

\section{Reproducibility}

The benchmark release is designed around command-line reproduction rather than
manual notebook execution. Every benchmark run writes a manifest containing the
benchmark version, command, Python version, source inputs, scenarios, seed
pairs, policies, cascade bands, label limit, evaluation load limit,
feasibility-layer version, output file paths, and row counts. Benchmark
submissions should include the manifest and all summary CSVs.

\section{Conclusion}

FreightBidBench provides a public, reproducible testbed for latency-aware
truckload bid acceptance. Its value is not that it solves the problem, but that
it gives researchers a shared environment where closed-loop profit, latency,
and operational feasibility can be compared under identical stochastic seeds.

\bibliographystyle{plainnat}
\bibliography{references}

\end{document}
